% Auto-generated by tools/qbe.py project-article-update.
% Do not edit this file by hand; edit the proof artifacts or article sections instead.

\section{Current Cycle Status}
\label{app:generated-cycle-status}

This appendix is intentionally short.  The full run logs, Chinese summaries,
and ChatGPT Pro prompts stay in the repository; the Overleaf report includes
only the compact status needed to keep the manuscript honest and fast to
compile.

\paragraph{Latest cycle.}
\texttt{QBE-AUTO-002}, cycle \texttt{1};
generated at \texttt{2026-06-17 02:10:36}.

\paragraph{Lean gate signal.}
\begin{itemize}
  \item \texttt{QuantumBlockEncoding/RobinMatrix.lean:26192:  sorry}
  \item \texttt{QuantumBlockEncoding/RobinMatrix.lean:26222:  sorry}
\end{itemize}

\paragraph{Plain-language status for this case study.}
The current GHL case study is not blocked because the paper lacks a proof sketch.  It is blocked because ABEIS must turn the circuit in the paper into an exact matrix statement and prove that the clean block entry has the coefficient claimed by the paper.  In ordinary terms, the paper says that a specific path through the circuit leaves the desired coefficient, while Lean requires the corresponding matrix entry and all unwanted branches to be proved exactly.

\paragraph{Current proof-DAG frontier.}
\begin{itemize}
  \item \texttt{projection\_interface\_normalizer\_bridge}: compiled route memory via \texttt{oneTermRobinGamma3BoundaryTheoremFacingProjectionInterface\_normalizerEval\_n3}; no theorem-facing semantic flag is promoted.
  \item \texttt{active\_prepared\_circuit\_field\_target}: compiled package \texttt{oneTermRobinGamma3BoundaryActivePreparedCircuitFieldTarget\_n3 H env}; it records the active seven-gate entry versus prepared singleton clean entry and keeps \texttt{activePreparedEntryEqualityProved = false}.
  \item \texttt{corrected\_finite\_block\_projection\_equality}: active audit frontier.  The Lean-facing field is \texttt{(oneTermRobinGamma3BoundarySourcePreparedProjectionTarget\_n3 H env).activeToPreparedSingletonEvalStatement}, equivalently \texttt{oneTermRobinGamma3BoundaryActivePreparedCompositeEvalStatement\_n3 H env} or \texttt{(oneTermRobinGamma3BoundaryActivePreparedEntryTarget\_n3 H).entryEqualityStatement}.
  \item \texttt{fixed\_product\_to\_coefficient}: blocked root \texttt{oneTermRobinGamma3ProductToCoefficientObligation 3 0 0}, waiting on the corrected finite block/projection equality and final coefficient bridge.
\end{itemize}

\paragraph{Open first-case-study obligations.}
\begin{itemize}
  \item \texttt{Uindic}: Indicator unitary; status Permutation and self-inverse helpers compile, but this is not yet the final block-encoding proof..
  \item \texttt{RyBoundary}: Boundary-controlled Ry rotations; status Active convention audit: the Ry angle convention must be source-supported before closing the boundary amplitude proof..
  \item \texttt{RobinTheorem}: One-term Robin block-encoding theorem; status Main active target: the theorem-facing Lean statement is not closed..
  \item \texttt{GammaSlices}: Gamma wavefunction slices and coefficient product; status Several finite boundary lemmas exist; the final projection/product bridge is still open..
\end{itemize}

\paragraph{Open cited-contract obligations.}
\begin{itemize}
  \item \texttt{tl-ghl-lemma1-banded-sparse-access}: contract-only; next action Keep as explicit external contract until the exact cited construction is formalized or imported as a theorem..
  \item \texttt{tl-ghl-lemma3-sparse-amplitude}: contract-only; next action Maintain the clean-branch contract and isolate any missing unitarity proof as an external technical lemma..
  \item \texttt{tl-ghl-theorem5-piecewise-polynomial-of}: contract-only; next action Use only as a named contract in the current theorem; do not invent stronger smoothness or range assumptions..
\end{itemize}

\paragraph{Recent verifier feedback.}
\begin{itemize}
  \item Leaf \texttt{theorem\_facing\_projection\_interface\_normalizer\_bridge}, class \texttt{stale\_leaf}; next route: retire normalizer bridge as accepted route memory; middle prepares next coefficient bridge; lower2 should not duplicate the existing theorem.
  \item Leaf \texttt{theorem\_facing\_projection\_interface\_normalizer\_bridge}, class \texttt{stale\_leaf}; next route: Retire this leaf for further lower2 assignment. Middle should update the DAG to the next narrow source-backed blocker after the interface-field normalizer bridge, without promoting product-to-coefficient, normalized-block, LCU, block, final-extraction, oracle, unitarity, or resource flags..
\end{itemize}

\paragraph{Manuscript rule.}
The project report may discuss this cycle as process evidence.  It must not
claim completion of the first case study \citep{guseynovHuangLiu2025} unless
the final theorem-facing block-extraction statement is closed in Lean and the
Markdown/LaTeX proof map has been synchronized.
